% =========================================================================
%  Ejercicio 2b - Ensayo individual sobre el articulo
%  Responsable: [nombre de cada integrante]
%
%  NOTA: CADA integrante del equipo debe escribir SU propa seccion. Dupliquen
%  este archivo (p.ej. ejercicio_2b_<apellido>.tex) o agregen subsecciones y
%  luego registren cada archivo en main.tex.
% =========================================================================

\section*{Ejercicio 2b. Ensayo sobre el artículo}

% Extension: entre 1 y 2 cuartillas por integrante.
%
% Estructura obligatoria (en parrafos, sin listas):
% - INTRODUCCION: presenta brevemente el tema, por que es importante, y
%   plantea tu postura o idea principal.
% - DESARROLLO: resume los argumentos del autor, relaciona con otros
%   conocimientos/lecturas/ejemplos, expone tus propios argumentos (a favor
%   o en contra) con razones, indica el objetivo que planteo el autor y como
%   se relaciona con la materia, y la tematica central y temas laterales y su
%   relacion con tu practica profesional.
% - CONCLUSION: reafirma tu postura, que aprendiste o la aportacion mas
%   valiosa, y si el tema abre preguntas o retos futuros.
% - FORMATO: parrafos completos, conectores logicos, ortografia y gramatica
%   cuidadas.

% Referencia del articulo (APA) - completar los datos:
% \cite{datavalue}